\documentclass[aspectratio=169,10pt,spanish]{beamer}

\input{../../../_preambulo_beamer.tex}
\input{../../../_comandos_beamer.tex}

\selectlanguage{spanish}

\title{MA1001B - Modelación Estadística para la Toma de Decisiones}
\subtitle{Lección 20: Distribución uniforme continua y Monte Carlo}
\author{}
\institute{Tecnol\'ogico de Monterrey}
\date{}

\begin{document}

\begin{frame}[plain]
  \vspace{1.2cm}
  {\Large\bfseries MA1001B - Modelación Estadística para la Toma de Decisiones\par}
  \vspace{0.7cm}
  {\large Lección 20: Uniforme continua y Monte Carlo\par}
  \vspace{0.7cm}
  {\color{TecRojo}\rule{\textwidth}{1.2pt}}
  \vspace{0.8cm}

  Facultad MA1001B\\[0.3cm]
  Periodo\\[0.3cm]
  Tecnol\'ogico de Monterrey
\end{frame}

\begin{frame}{La pregunta de densidad plana}
  \begin{alertblock}{Pregunta guía}
    Si el tiempo de procesamiento es igualmente probable en cualquier punto de
    $[8,20]$ minutos, ¿cómo se calcula $P(X>16)$ y cómo puede la simulación
    comprobarlo?
  \end{alertblock}

  \vspace{0.6em}
  Al final de esta lección, usted puede:
  \begin{itemize}
    \item escribir la densidad plana $f(x)=1/(b-a)$;
    \item calcular una cola como longitud restante sobre $b-a$;
    \item comprobar esa cola con una muestra Monte Carlo sembrada;
    \item explicar por qué una densidad acumulada rompe el modelo uniforme.
  \end{itemize}

  \vfill
  \scriptsize Todos los tiempos de procesamiento usados hoy son sintéticos. No se usan datos reales de empresa.
\end{frame}

\begin{frame}{Minutos de procesamiento Uniforme$(8,20)$}
  \small
  Altura de densidad y media:
  \[
    f(x)=\frac{1}{20-8}=\frac{1}{12}=0.083333,
    \qquad
    E[X]=\frac{8+20}{2}=14.
  \]
  Cola por longitud:
  \[
    P(X>16)=\frac{20-16}{12}=\frac{1}{3}=0.333333.
  \]

  \vspace{0.3em}
  \begin{exampleblock}{Reloj sintético de tickets}
    Cada minuto entre 8 y 20 es igualmente probable. La probabilidad es
    longitud, no una altura puntual.
  \end{exampleblock}
\end{frame}

\begin{frame}[fragile]{Paso 1: Cola uniforme exacta}
  \begin{columns}[T]
    \begin{column}{0.42\textwidth}
      \small
      \begin{lstlisting}
model = uniform(
  loc=8, scale=12
)
p_gt_16 = 1 - model.cdf(16)
      \end{lstlisting}

      \begin{block}{Salida verificada}
        Altura $=0.083333$\\
        Media $=14.000000$\\
        Razón de longitudes $=0.333333$\\
        $P(X>16)=0.333333$
      \end{block}
    \end{column}
    \begin{column}{0.58\textwidth}
      \centering
      \includegraphics[width=\textwidth]{../figuras/uniforme_mc_01_uniforme_exacta.png}
      \vspace{0.2em}
      \scriptsize Densidad plana y cola sombreada del Paso 1
    \end{column}
  \end{columns}
\end{frame}

\begin{frame}{Monte Carlo recupera la probabilidad por longitud}
  \small
  Extraiga $n=100{,}000$ veces de $\mathrm{Uniforme}(8,20)$ con semilla 42.
  Cuente la fracción de extracciones por encima de 16:
  \[
    \widehat{P}(X>16)=\frac{\#\{X_i>16\}}{100{,}000}.
  \]
  La cola simulada es $0.333290$. El error absoluto frente a $1/3$ es
  $0.000043$. La media simulada es $14.007499$ frente a la exacta $14$.
\end{frame}

\begin{frame}[fragile]{Paso 2: $100{,}000$ extracciones con semilla 42}
  \begin{columns}[T]
    \begin{column}{0.40\textwidth}
      \small
      \begin{lstlisting}
draws = model.rvs(
  size=100_000,
  random_state=rng,
)
p_hat = np.mean(draws > 16)
      \end{lstlisting}

      \begin{block}{Salida verificada}
        Cola exacta $=0.333333$\\
        Monte Carlo $=0.333290$\\
        Error abs.\ $=0.000043$\\
        Media MC $=14.007499$
      \end{block}
    \end{column}
    \begin{column}{0.60\textwidth}
      \centering
      \includegraphics[width=\textwidth]{../figuras/uniforme_mc_02_monte_carlo.png}
      \vspace{0.2em}
      \scriptsize Histograma frente a densidad exacta del Paso 2
    \end{column}
  \end{columns}
\end{frame}

\begin{frame}{Paso 3: La uniforme falla cuando la densidad se acumula}
  \begin{columns}[T]
    \begin{column}{0.42\textwidth}
      \small
      Uniforme $P(X>16)=0.333333$.

      \vspace{0.4em}
      Acumulación triangular derecha en 20:
      \[
        P(X>16)=0.555556.
      \]
      Media acumulada $=16$, no $14$.

      \vspace{0.4em}
      \textbf{Límite:} un modelo plano es incorrecto cuando la densidad se
      acumula en un extremo. La Lección 21 estandariza un reloj en forma de
      campana.
    \end{column}
    \begin{column}{0.58\textwidth}
      \centering
      \includegraphics[width=\textwidth]{../figuras/uniforme_mc_03_limite_densidad_acumulada.png}
      \vspace{0.2em}
      \scriptsize Densidades plana y acumulada del Paso 3
    \end{column}
  \end{columns}
\end{frame}

\begin{frame}{Verificación de conceptos}
  \begin{enumerate}
    \item ¿Por qué $f(x)=1/12$ en $[8,20]$ es una densidad, no una probabilidad?
    \item Calcule $P(X>16)$ a partir de la longitud restante $4$ y la longitud total $12$.
    \item La cola Monte Carlo es $0.333290$. ¿Sustituye eso a $1/3$?
    \item ¿Por qué $0.555556$ es el número equivocado para un reloj uniforme?
  \end{enumerate}

  \vspace{0.6em}
  \begin{block}{Conexión con la Lección 21}
    La sesión puede continuar directamente con la curva normal estándar y las
    puntuaciones $z$.
  \end{block}

  \vfill
  \scriptsize Fuente: OpenStax, \textit{Introductory Business Statistics 2e},
  Capítulo 5, Sección 5.2. CC BY-NC-SA.
\end{frame}

\appendix



\end{document}
